% !TEX program = lualatex
% glossa-template.tex
% Copyright 2016 Guido Vanden Wyngaerd
%
% This work may be distributed and/or modified under the
% conditions of the LaTeX Project Public License.
% The latest version of this license is in
%   http://www.latex-project.org/lppl.txt
% and version 1.3 or later is part of all distributions of LaTeX
% version 2005/12/01 or later.
%
% This work has the LPPL maintenance status `maintained'.
% 
% The Current Maintainer of this work is 
% Guido Vanden Wyngaerd (guido.vandenwyngaerd@kuleuven.be).
%
% This work consists of the files 
% glossa.cls
% glossa.bst
% gl-authoryear-comp.cbx
% biblatex-gl.bbx
% glossa-template.tex
% glossa.png
%
% The files of the work are derived from the Semantics & Pragmatics style files
% by Kai von Fintel, Christopher Potts, and Chung-chieh Shan
% All changes are documented on the github repository 
% https://github.com/guidovw/Glossalatex.

\PassOptionsToPackage{table}{xcolor}
\PassOptionsToPackage{xcolor}{dvipsnames}
\documentclass[times,linguex]{glossa}
\usepackage{rotating}
\usepackage{tablefootnote}
\usepackage{colortbl}
\usepackage{color}
\usepackage{multicol}
\usepackage{booktabs}

\usepackage{adjustbox}
\usepackage{array}

\newcolumntype{R}[2]{%
>{\adjustbox{angle=#1,lap=\width-(#2)}\bgroup}%
l%
<{\egroup}%
}

\newcommand*\rots{\multicolumn{1}{R{90}{.7em}}}% no optional argument here, please!
%\usepackage{xcolor}

% possible options:
% [times] for Times font (default if no option is chosen)
% [cm] for Computer Modern font
% [lucida] for Lucida font (not freely available)
% [brill] open type font, freely downloadable for non-commercial use from http://www.brill.com/about/brill-fonts; requires xetex
% [charis] for CharisSIL font, freely downloadable from http://software.sil.org/charis/
% for the Brill an CharisSIL fonts, you have to use the XeLatex typesetting engine (not pdfLatex)
% for headings, tables, captions, etc., Fira Sans is used: https://www.fontsquirrel.com/fonts/fira-sans
% [biblatex] for using biblatex (the default is natbib, do not load the natbib package in this file, it is loaded automatically via the document class glossa.cls)
% [linguex] loads the linguex example package
% !! a note on the use of linguex: in glossed examples, the third line of the example (the translation) needs to be prefixed with \glt. This is to allow a first line with the name of the language and the source of the example. See example (2) in the text for an illustration.
% !! a note on the use of bibtex: for PhD dissertations to typeset correctly in the references list, the Address field needs to contain the city (for US cities in the format "Santa Cruz, CA")

%\addbibresource{sample.bib}
% the above line is for use with biblatex
% replace this by the name of your bib-file (extension .bib is required)
% comment out if you use natbib/bibtex

\let\B\relax %to resolve a conflict in the definition of these commands between xyling and xunicode (the latter called by fontspec, called by charis)
\let\T\relax
\usepackage{xyling} %for trees; the use of xyling with the CharisSIL font produces poor results in the branches. This problem does not arise with the packages qtree or forest.
%\usepackage[linguistics]{forest} %for nice trees!




% \pdf* commands provide metadata for the PDF output. ASCII characters only!
\pdfauthor{}
\pdftitle{Comparing at-issueness diagnostics}
\pdfkeywords{}

\title[Comparing at-issueness diagnostics]{An experimental comparison of at-issueness diagnostics\\ 
\bigskip \large
Word count (including references, excluding abstract and supplements): 11,352
}
% Optional short title inside square brackets, for the running headers.

\author[]% short form of the author names for the running header. If no short author is given, no authors print in the headers.
{%as many authors as you like, each separated by \AND.
% \spauthor{Waltraud Paul\\
% \institute{CRLAO, CNRS-EHESS-INALCO}\\
% \small{%105, Bd. Raspail, 75005 Paris\\
% waltraud.paul@ehess.fr}
% }
% \AND
% \spauthor{Guido Vanden Wyngaerd \\
% \institute{KU Leuven}\\
% \small{%Warmoesberg 26, 1000 Brussel\\
% guido.vandenwyngaerd@kuleuven.be}
% }%
}

\input{author-added}

% positive coefficients/difference
\definecolor{red}{RGB}{178,24,43}

% negative coefficients/difference
\definecolor{blue}{RGB}{33,102,172}

% comments by JT
\newcommand{\jt}[1]{%TC:ignore
\textbf{\color{blue}JT: #1}%TC:endignore
}
\newcommand{\lh}[1]{%TC:ignore
\textbf{\color{orange}LH: #1}%TC:endignore
}

\begin{document}



\section{Experiments 1-4 \label{sec:2_experiments}}

To try to understand the greater by-content differentiation of the `asking whether' diagnostic, we focus on properties that set it apart from the other three diagnostics. Two differences are particularly salient: First, it is the only diagnostic where the target content itself occurs in a polar question, whereas in the other diagnostics, the target content appears in a declarative sentence. Second, it is the only diagnostic where participants are asked directly about the intentions of the speaker, while the other ones collect meta-linguistic judgments like naturalness ratings or continuation preferences. To distinguish between these possibilities, the next section compares the `asking whether' diagnostic (Exp.~5) to a diagnostic that also embeds the target content in a polar question but elicits naturalness ratings on direct responses (Exp.~6). 




\subsection*{Why do the different experiments/diagnostics give different results? \label{ssec:discussion-questions}}

\bigskip
\noindent
Possible explanations:

\begin{enumerate}
  
  \item For the CCs of the five clause-embedding predicates: Exp 2 shows great differentiation (know $<$ confess, discover $<$ confirm $<$ be right; 9/10 cells), Exp 3 shows no differentiation (0/10 cells, and Exps 1 (6/10 cells, but be right wonkiness) and 4 (5/10) show much less differentiation
  
  \begin{itemize}
    
    \item {\bf Hypothesis:} The greater differentiation in Exp 2 is because of the embedding polar questions, not because of the response task.
    
    \item We confirmed this by running Exps 5-6, where we found lots of differentiation among the CCs of the 20 clause-embedding predicates.
    
    \item Why would polar question-embedding produce greater differentiation than declarative embedding in the five predicates featured in our experiments?
    
    \item Observation: We picked the five predicates because they differed greatly in the asking-whether diagnostic. But the CCs of all five predicates received pretty high ratings in gradient and categorical inference experiments of Degen/Tonhauser/Language, so there's a pretty strong inference that the CC follows from these utterances (not all entailed; confess $<$ know $<$ confirm $<$ discover $<$ be right); only one that is entailed is be right and know. 
    
    \item So the five predicates don't differ a great deal in declaratives when it comes to speaker commitment (? or interpreter inference) to the CC
    
    \item So when a speaker in Exps 1, 3-4 asserts the declarative, the interpreter is likely to take both contents to follow. These three diagnostics may not reliably distinguish between at-issueness but are rather about availability for subsequent reference. And since all of the CCs are (almost) entailed, they are available for subsequent reference.
    
    \item Exp 2 is the only diagnostic where MC and CC are both embedded under entailment-canceling operator; so here we entailments are canceled and the diagnostic can pay more attention to at-issueness? the more projective the more not-at-issue
    
    \item \citealt{scontras-tonhauser2025,scontras-tonhauser2026}: q-know utterance has more utility when CC is assumed to be true than when BEL is assumed to be true. ==> explains low at-issueness rating for CC in Exp 2
    
  \end{itemize}
  
\end{enumerate}





\bibliography{at-issueness} %for use with natbib; comment out if you use biblatex, and change 'sample' by the name of your bib-file

%TC:endignore
\end{document}
%%% Local Variables:
%%% mode: latex
%%% TeX-master: t
%%% TeX-engine: luatex
%%% End:
